
\section{Abstração Metalinguística}

\begin{frame}
  \frametitle{De Programadores a Designers de Linguagens}
  \excerptpage{It is no exaggeration to regard this as the most fundamental idea in programming: The evaluator, which determines the meaning of expressions in a programming language, is just another program. To appreciate this point is to change our images of ourselves as programmers. We come to see ourselves as designers of languages, rather than only users of languages designed by others. In fact, we can regard almost any program as the evaluator for some language.}{Gerald Jay Sussman \& Harold Abelson}
\end{frame}

\begin{frame}
  \frametitle{Por que um Interpretador?}
  \begin{itemize}
    \item Desmistificação do interpretador; \\ \vspace{1cm}
    \pause
    \item Flexibilidade e extensão; \\ \vspace{1cm}
    \pause
    \item Criação de \textit{Domain Specific Languages} (DSLs).
  \end{itemize}
\end{frame}

\begin{frame}
  \excerptpage{Once you have the interpreter in your hands, you have all this power to start playing with the language. [...] There's this notion of metalinguistic abstraction, which says [...] that you can gain control of complexity by inventing new languages, sometimes. One way to think about computer programming is that it only incidentally has to do with getting a computer to do something. Primarily, what a computer program has to do with is a way of expressing ideas, of communicating ideas. Sometimes, when you want to communicate new kinds of ideas, you'd like to invent new modes of expressing them.}{Harold Abelson}
\end{frame}